\documentclass[11pt]{article}
\usepackage[margin=1in]{geometry}
\usepackage{amsmath,amssymb,amsthm,booktabs,hyperref}
\newtheorem{proposition}{Proposition}
\newtheorem{remark}{Remark}
\newcommand{\R}{\mathbb R}
\title{An analytic rationale for an anchored two-weight relaxation baseline}
\author{}
\date{September 27, 2026}
\begin{document}
\maketitle
\noindent
This note derives a small, explicitly delimited baseline-design result. It
does not prove that the observed pair $(2.5,1)$ minimizes actual BoomerAMG
runtime, nor that a full multilevel V-cycle equals a single Richardson step.
The mathematical claims are separate from the Module 05 Run 04 measurements. The weighted polynomial-smoothing
objective is classical; the anchored specialization and two-grid spectral
calculation below are provided with self-contained proofs, without a claim
of priority.

\section{Why a difference of weights is not an effective relaxation}
For $S(w)=I-wT$, two successive steps produce
\[
 S(b)S(a)=I-(a+b)T+abT^2.
\]
With $m=(a+b)/2$ and $d=(a-b)/2$, the scalar polynomial is
\[
 p_{a,b}(t)=(1-mt)^2-d^2t^2.
\]
There is no general rule identifying $a-b$ with an optimal constant weight.
Even the usual first-kind Chebyshev construction, on an SPD interval
$[\mu,L]$, instead has
\[
 a^{-1}+b^{-1}=L+\mu,\qquad
 w_{\rm constant}=\frac{2}{L+\mu}=\frac{2ab}{a+b}.
\]
This relation holds for those interval-minimax choices, not arbitrary
runtime-optimal AMG weights. For $(a,b)=(2.5,1)$ the latter harmonic mean
is $10/7$, not $1.5$ or $1.6$.

\section{The normalized SPD smoothing model}
For symmetric positive definite $A$, define the conventional row-$\ell_1$
Jacobi diagonal
\[
 D_{ii}=a_{ii}+\sum_{j\ne i}|a_{ij}|.
\]
For every vector $x$,
\[
 x^\top Ax
 \le \sum_i a_{ii}x_i^2+
       \sum_{i<j}|a_{ij}|(x_i^2+x_j^2)
 =x^\top Dx.
\]
Consequently $0\prec A\preceq D$. In energy coordinates the smoother is
\[
 S(w)=I-wT,\qquad T=A^{1/2}D^{-1}A^{1/2},\qquad0\prec T\preceq I.
\]
The equivalent symmetric matrix $D^{-1/2}AD^{-1/2}$ has the same
eigenvalues. The row-$\ell_1$ formulation and polynomial acceleration are
discussed explicitly by D'Ambra et al.~\cite{dambra}.

For polynomial smoothers, a classical two-grid bound contains
$\sup_{0<t\le1}\sqrt t\,|p(t)|$; see Lottes~\cite{lottes}, Eq.~(4).
We therefore consider its square, $\sup_{0\le t\le1}t p(t)^2$.
This is a smoothing-bound surrogate, not a wall-clock objective.

\section{An anchored two-stage minimax problem}
Retain one ordinary step of weight one and allow a second weight $a$.
The paired polynomial is
\[
 p_a(t)=(1-t)(1-at),\qquad
 \eta(a)=\max_{0\le t\le1}t(1-t)^2(1-at)^2.
\]

\begin{proposition}[Anchored weighted-minimax coefficient]
The unique minimizer of $\eta(a)$ over $a\in\R$ (hence also over
$a\in[1,3]$) is
\[
 a_*=\frac{3+\sqrt5}{2}=2.6180339887\ldots,
 \qquad
 \eta(a_*)=\frac{10-2\sqrt5}{125}.
\]
\end{proposition}
\begin{proof}
Set $f_a(t)=\sqrt t(1-t)(1-at)$. For $0<t<1$,
\[
 f_a'(t)=\frac{1-3(a+1)t+5at^2}{2\sqrt t}.
\]
At $a=a_*$ its only interior stationary points are
\[
 t_+=1-\frac2{\sqrt5},\qquad
 t_-=\frac12\left(1+\frac1{\sqrt5}\right).
\]
Direct substitution gives
\[
 f_{a_*}(t_+)=g,\qquad f_{a_*}(t_-)=-g,\qquad
 g^2=\frac{10-2\sqrt5}{125}.
\]
Both endpoint values are zero, so $\max_t|f_{a_*}(t)|=g$.
If $a<a_*$, then
\[
 f_a(t_+)=g+(a_*-a)t_+^{3/2}(1-t_+)>g.
\]
If $a>a_*$, then
\[
 f_a(t_-)=-g-(a-a_*)t_-^{3/2}(1-t_-)<-g.
\]
Every other $a$ therefore has a larger maximum absolute value. Squaring
proves the result.
\end{proof}

For the five-point grid used in the development search, the exact
stationary-point calculation gives:
\begin{center}
\begin{tabular}{rr}
\toprule
$a$ & $\eta(a)$\\
\midrule
1.0 & 0.0819200000\\
1.5 & 0.0652127357\\
2.0 & 0.0538822475\\
2.5 & 0.0457988174\\
3.0 & 0.0762755014\\
\bottomrule
\end{tabular}
\end{center}
Thus $(2.5,1)$ is the coarse-grid minimizer of this independently specified
surrogate. On the actual $0.05$ action grid, the minimizer is $(2.6,1)$,
with $\eta(2.6)=0.0444568225$. This discrete statement follows by the
same calculation; it is not justified solely by nearest rounding.
The coarse-grid objective is about $3.56\%$ above the continuous minimum.
This percentage is a surrogate-gap percentage, not a predicted runtime gap.

\begin{remark}
Although $a_*-1=(1+\sqrt5)/2$ is near $1.6$, it is not the optimum
one-step weight for the same weighted smoothing problem. That one-step
optimum is $4/3$~\cite{lottes}. The proximity must not be turned into
an explanation of the experimentally observed best fixed AMG weight.
\end{remark}

\section{A precise connection to complete two-grid cycles}
The preceding objective is not only a scalar heuristic: it also bounds a
particular complete two-grid periodic iteration.

\begin{proposition}[Compression identity and two-cycle rate bound]
Let $0\prec T=T^\top\preceq I$ and let $Q=Q^\top=Q^2$ be the exact
coarse-correction projection in energy coordinates. Define
\[
 E(w)=(I-wT)Q(I-wT),\qquad P_a=p_a(T),\qquad
 K=\|T^{-1/2}Q\|_2^2.
\]
Then
\[
 \rho\big(E(1)E(a)\big)
 =\|QP_aQ\|_2^2
 \le K\,\eta(a).
\]
The same spectral radius is obtained with the opposite phase
$E(a)E(1)$.
\end{proposition}
\begin{proof}
Write $S_1=I-T$ and $S_a=I-aT$. These commute, so
\[
 E(1)E(a)=S_1 QP_a Q S_a.
\]
Cyclic permutation of square factors preserves the characteristic
polynomial, giving the same eigenvalues as
\[
 S_aS_1QP_aQ=(P_aQ)^2.
\]
In a basis adapted to $\operatorname{range}Q$ and $\ker Q$, $P_aQ$
has the eigenvalues of the symmetric compression $QP_aQ$ on
$\operatorname{range}Q$, together with zeros. Therefore
\[
 \rho\big(E(1)E(a)\big)=\rho(P_aQ)^2=\|QP_aQ\|_2^2.
\]
Moreover,
\[
 \|QP_aQ\|_2
 \le\|P_aQ\|_2
 \le\|P_aT^{1/2}\|_2\,\|T^{-1/2}Q\|_2
 \le\sqrt{\eta(a)K}.
\]
The two products $E(1)E(a)$ and $E(a)E(1)$ have the same eigenvalues.
\end{proof}

This minimizes an upper bound for a two-cycle spectral radius with fixed
$K$, not the exact rate of every hierarchy. The bound can exceed one
when coarse approximation is poor. It does not control intermediate
growth, finite-tolerance stopping time, or unequal execution costs.
In a genuine multilevel cycle, changing relaxation on all levels also
changes the inexact coarse correction. Its effective correction is
not necessarily this fixed projection $Q$. Nonsymmetric advection
requires a separate analysis.

In the manuscript's rank-one special case, $Q=uu^\top$,
\[
 QP_aQ=
 \big[1-m_1(1+a)+m_2a\big]Q,\quad
 m_1=u^\top Tu,\quad m_2=u^\top T^2u.
\]
The exact cancellation test for $(2.5,1)$ is therefore
$1-3.5m_1+2.5m_2=0$. Knowing $w_*=m_1/m_2$ alone does not determine it.

\section{An unanchored literature comparator}
The degree-two fourth-kind Chebyshev polynomial in Lottes~\cite{lottes},
Eq.~(7), is
\[
 p_2(t)=1-4t+\frac{16}{5}t^2
       =(1-at)(1-bt),
\quad
 a=2+\frac2{\sqrt5},\quad b=2-\frac2{\sqrt5}.
\]
It gives $(a,b)\simeq(2.894427,1.105573)$, or $(2.90,1.10)$ on the
$0.05$ grid. Exact factorization realizes this polynomial with two
successive Jacobi steps on the same matrix; rounding changes the
polynomial. Using these coefficients on successive complete multilevel
V-cycles is a motivated transfer, not the original polynomial-smoother
theorem. Standard SRJ/nonstationary Richardson theory also concerns
operator products rather than differences of coefficients
\cite{yang,adsuara}.

The sharper multilevel objective in Lottes, Section 3.2, is different:
$\sup_t t p(t)^2/(1-p(t)^2)$. Its unrestricted degree-two optimum
has weights $\sqrt5/2$ and $(5+\sqrt5)/2$, the latter exceeding the
study's allowed upper weight 3. We do not call the anchored minimizer
optimal for that different objective, or clip another optimum and
retain its theorem.

\section{Audit against the implemented experiment}
The production profile is \texttt{l1\_jacobi\_direct\_coarse}. The native
source sets relaxation type 18 on the down and up stages and type 9
(Gaussian elimination) on the coarsest stage. Natural relaxation order
(\texttt{relax\_order=0}) is the HYPRE default and is not overridden in
any of the 600 frozen test jobs. For this order, HYPRE constructs type-18
norms with option 1 and a null C/F marker, giving exactly
$D_{ii}=\sum_j|a_{ij}|$. Since an SPD matrix has $a_{ii}>0$, this equals
the diagonal used above. The type-18 kernel uses the old iterate on all
points and applies $x\leftarrow x+wD^{-1}(b-Ax)$.

The inputs are the positive-coefficient, zero-advection 7-point diffusion
stencil with homogeneous Dirichlet boundaries; hence the fine matrix is
SPD. Standard Galerkin coarsening preserves SPD in exact arithmetic when
interpolation has full column rank. Positive off-diagonal entries on a
coarse level would not invalidate the absolute-row-sum domination proof.
This note does not certify every computed coarse matrix against
floating-point perturbations.

Each experimental action runs one complete V-cycle, with one pre-sweep
and one post-sweep. The selected weight is set on every relaxation level;
the direct coarse solve is unaffected by that weight. Thus the actual
fine-level error operator has the form
\[
 E_0(w)=S_0(w)\,[I-P_0 B_1(w)P_0^\top A_0],S_0(w)
\]
in original coordinates, where here
$S_0(w)=I-wD_0^{-1}A_0$ and $B_1(w)$ is the recursive coarse solver.
For more than two levels, $B_1(w)$ generally depends on $w$ and is not
$A_1^{-1}$. The fixed-projection proposition therefore applies exactly
to the ideal two-grid specialization, not to arbitrary production cycles.
The coefficient is a polynomial-smoothing-inspired baseline.

\section{Run 04 protocol and paper wording}
Module 05 Run 04 repeats the Run 03 diffusion $60^3$ evaluation with the
same six frozen setup/controller checkpoints (including the refreshed
seed 4), the same 100 test inputs, identical setup tuples, three timing
repetitions, and the same previously selected fixed weights. The only
policy replacement is $(2.5,1)\mapsto(2.6,1)$. The phase is prescribed as
high first and reset on each primary solve; the older $(1,3)$ comparator
remains. This phase is a protocol choice, not a minimax theorem.

The tolerance remains $10^{-6}$ relative to the initial residual, with
the existing 50-cycle stopping semantics and charged reference recovery.
Native continuation costs include all primary solves and recovery
setup/solve; the common initial setup is excluded. Controller-inclusive
costs are reported separately. No learner is retrained, no fixed-weight
scan is repeated, and no schedule is selected using Run 04 timings.
The historical development use and checkpoint selection remain disclosed;
this is not a new independent replication.

A paper-ready description is:
\begin{quote}
We include a prescribed period-two relaxation baseline with weights
$(2.6,1)$, motivated by normalized SPD $\ell_1$-Jacobi polynomial smoothing.
Anchoring one coefficient at unity and minimizing
$\max_{t\in[0,1]}t(1-t)^2(1-at)^2$ gives
$a_*=(3+\sqrt5)/2$; analytic evaluation on our $0.05$ grid selects $a=2.6$.
An exact two-grid model links this surrogate to an upper bound on the
two-cycle spectral radius. We transfer the coefficients to complete
multilevel cycles and assess their finite-tolerance cost empirically.
\end{quote}

The accompanying verification script checks the symbolic extrema,
all 41 grid coefficients, 500 small-matrix compression identities and
bounds, every frozen test input, and one native profile per checkpoint.
These checks supplement, rather than replace, the proofs.

\begin{thebibliography}{9}
\bibitem{yang}
X. I. A. Yang and R. Mittal,
Acceleration of the Jacobi iterative method by factors exceeding 100
using scheduled relaxation,
\emph{Journal of Computational Physics}, 274 (2014), 695--708.
doi:10.1016/j.jcp.2014.06.010.
\bibitem{adsuara}
J. E. Adsuara, I. Cordero-Carri\'on, P. Cerd\'a-Dur\'an, V. Mewes,
and M. A. Aloy,
On the equivalence between the Scheduled Relaxation Jacobi method
and Richardson's non-stationary method, arXiv:1607.03712 (2016).
\bibitem{lottes}
J. Lottes, Optimal polynomial smoothers for multigrid V-cycles,
\emph{Numerical Linear Algebra with Applications}, 30(6) (2023), e2518.
doi:10.1002/nla.2518. Equation numbers here follow arXiv:2202.08830v3.
\bibitem{dambra}
P. D'Ambra, F. Durastante, S. Filippone, S. Massei, and S. Thomas,
Optimal Polynomial Smoothers for Parallel AMG,
\emph{Numerical Algorithms} (2025).
doi:10.1007/s11075-025-02117-6; arXiv:2407.09848v3.
\end{thebibliography}
\end{document}
